% Transcription — fonds Alexandre Grothendieck, shelfmark 137, batch 4 (pages 61-80).
% Established from the facsimile archives/batches/137/batch-04.pdf (a copy of
% 137.pdf fetched through the project's relay,
% https://grothendieck-relay.onrender.com/source/137.pdf; PDF page k is
% archivists' page 60+k, no cover sheet), per the skill transcribe-grothendieck.
% The folder is eighty-five pages in five batches (1-20, 21-40, 41-60, 61-80,
% 81-85), transcribed in parallel. Batches 3 and 5 were read before writing
% and this batch is aligned on them: pages 61-64 continue the run
% « Construction canonique de P (k-lin.) » of pages 57-60 (notation K, C^•,
% Psi^i, Phi^•, P_0, P_1, Sigma as there), and pages 81-85 continue page 80.
%
% Pass: Opus 5.5 (claude-opus-5-5), 2026-09-26 — first pass, unchecked against the pages by a human.
%
% Pages skipped: none. Pages summarised: none. All twenty pages are his
% mathematics and are transcribed. Page 73 carries only two lines (the rest
% of the sheet is blank, with show-through from page 72).
%
% His pagination: none on these leaves. Numbered runs and labels: the
% « Remarque » of p. 61; « Ex. » (p. 71) and « Ex 2 » (p. 72), the
% applications of multiplicative structures; a new run headed
% « Catégories polynomiales » (p. 74, underlined, set as a \subsection*)
% with axioms a), b), c) (p. 74), d) (p. 75), e), f), g) (p. 76) — the
% « axiomes a) à f) » cited on p. 83 are these, though the page goes on to
% g); the decomposition labelled (*) on p. 75; a second heading « Relations
% avec produits ⊗ et Γ^p » (p. 77, underlined, \subsection*); on p. 80 the
% maps (*) Γ^{pq} -> Γ^p Γ^q and the associativity square (*), which p. 81
% (batch 5) cites as « la condition d'associativité ci-dessus » before adding
% the unit conditions (**), (***). Recorded in \note{}s in place.
%
% Hand and apparatus. Blue ink throughout, fast; pp. 61, 69 and 74-75 carry
% long boxed-and-struck passages (given as \struck, with \ill{} where the
% cancelled words do not read); the struck lemma of pp. 75-76 is legible and
% is given unmarked with one \note saying it is cancelled. Pp. 70, 72 and
% 78 are the hardest: connective prose largely lost, formulas carry. Margin
% notes written sideways on pp. 64, 67, 70, 75, 76, 79, 80 are \marginal,
% mostly \ill{}. Tally: 223 \ill{} and 158 \uncertain{}.
%
% Notation fixed once for the batch, following batch 3: \mathbb{K} (sums of
% the universal complex), \mathbb{C}^• for the universal complex (his C is
% sometimes single-, sometimes double-struck; generic complexes of B stay
% C^•), \Psi^i, \Phi^•, \mathcal{P}, \mathcal{P}_0, \mathcal{P}_1 = \mathcal{Q}
% (his script Q), \Sigma; lower * for semi-simplicial realisations
% (\mathcal{P}_*, \mathbb{K}_*); * for his star product of cochain objects;
% from p. 74 on, \mathcal{P}, \mathcal{P}_1, \mathcal{P}_n(X,Y) and
% \mathbf{1} for his double-stroked unit, as in batch 5.
%
% What the batch is about. Pp. 61-64 end the characterisation of the
% category P of perfect ½-simplicial complexes: kernels and Σ-monos
% (Remarque, p. 61), the universal property of (P, C^•, Σ), work in
% P_1 = Q (complexes with L_i, Z_i, B_i, H_i free of finite type) where
% there are enough Σ-injectives (sums of C^•) and Σ-projectives, derived
% functors R^iF(Ψ^0) = H^i(F) and RF(Ψ^0) ≅ F(C^•) (pp. 62-64). From p. 64
% the multiplicative structures: the componentwise tensor product on the
% semi-simplicial realisation P_*, K_* an « ideal », P_0 and Q stable
% (Ψ^i ⊗ Ψ^j ≅ Ψ^{i+j} + an element of K); multiplicative structures on a
% functor F between ⊗-categories with the three compatibility diagrams
% (p. 65); for F : K -> B, i.e. a complex C^• of B, a « demi-simpliciale »
% multiplicative structure as maps (C^i ⊗ C^j) ⊗ Λ^iΦ ⊗ Λ^jΦ -> DP(C^•)
% (pp. 66-68); extension from K to Σ-left-exact functors on Q and P_0 and
% the unit problem (pp. 68-70); applications: a homotopy-associative,
% homotopy-commutative product C^• ⊗ C^• -> C^• making H^*(C^•) a graded
% ⊗-algebra (p. 71), the divided-power De Rham complex and a
% De Rham-Sullivan complex in characteristic 0 (pp. 72-73). Pp. 74-80 open
% « Catégories polynomiales »: a category P with a non-full subcategory
% P_1, axioms a)-g) (finite products, additivity of P_1, decomposition
% P(X_1 × ... × X_r, Y) = ⊕ P_n(X_1,...,X_r; Y) by homogeneity degree, a
% uniqueness lemma, degree of composites, symmetry, constants), k-polynomial
% categories; then ⊗ and Γ^n as representing objects, strict tensor
% products, Γ^n(X_1 × ... × X_r) ≅ ⊕ Γ^{ν_1}X_1 ⊗ ... ⊗ Γ^{ν_r}X_r,
% X ⊗ 1 ≅ X, and on p. 80 the functor Γ^* : P_1 -> Grad(P_1) turning sums
% into ⊗, with maps Γ^{pq} -> Γ^p ∘ Γ^q and their associativity square.
%
% Readings to check first: the first lines of p. 61 and its struck block;
% « coquilles » on p. 66; the prefix read \mathbb{F}^{ij} on p. 67; the
% long sentence of p. 72; the bottom of p. 70 and its vertical margin note;
% « trimultiplicatif » on p. 77; the struck interlinear lines of p. 78; the
% exponent \mathfrak{g} of P on p. 79; the superscripts in the square of
% p. 80.

\documentclass[11pt,a4paper]{article}
% Le préambule commun à toutes les transcriptions du fonds Grothendieck.
%
% Il tient en une idée : **l'incertitude se compose, elle ne se résout pas.**
% Une transcription qui lisse un mot douteux a détruit la seule chose pour
% laquelle on la faisait — un lecteur doit pouvoir savoir, à chaque endroit,
% ce qui était sur le feuillet et ce qui a été ajouté. Les six macros qui
% suivent sont donc l'appareil critique, et non de la décoration.
%
% Les mêmes commandes sont reconnues par `scripts/render.mjs`, qui produit la
% vue de lecture du site. Une macro ajoutée ici doit l'être aussi là-bas,
% faute de quoi le rendu échouera bruyamment — ce qui est voulu : un rendu
% silencieusement faux est pire qu'un rendu absent.

\ProvidesPackage{grothendieck}[2026/08/08 Transcriptions du fonds Grothendieck]

\usepackage{amsmath,amssymb,amsthm}
\usepackage{mathrsfs}
\usepackage{stmaryrd}
% Les diagrammes commutatifs sont partout dans ce fonds : c'est la langue
% même de la géométrie algébrique de Grothendieck, et une transcription qui
% les rendrait en prose ne transcrirait rien.
\usepackage{tikz-cd}
% Français par défaut — c'est la langue des feuillets — mais l'anglais est
% chargé aussi : la traduction et le résumé sont des documents anglais, et la
% typographie française (espaces avant les deux-points) y serait une faute.
% Ces documents basculent par \selectlanguage{english} en tête de corps.
\usepackage[english,french]{babel}
\usepackage{fontspec}
\usepackage{geometry}
\geometry{a4paper,margin=25mm,marginparwidth=28mm}
\usepackage{marginnote}
\usepackage{xcolor}
% ulem, not soul. `soul`'s \st and \ul are built on a character-by-character
% scan that XeLaTeX's font machinery breaks: the document fails with an
% undefined control sequence rather than degrading. ulem does the same two jobs
% and is Unicode-engine safe. `normalem` stops it hijacking \emph, which the
% transcriptions use for Grothendieck's own emphasis.
\usepackage[normalem]{ulem}
\usepackage{hyperref}
\hypersetup{colorlinks=true,linkcolor=black,urlcolor=black,pdfborder={0 0 0}}

% --- Métadonnées du lot -----------------------------------------------------
% Reprises telles quelles par le script de rendu, qui les affiche en tête de
% la vue de lecture. Elles fixent ce que le fichier prétend couvrir : sans
% elles, une transcription flotte sans point d'attache dans un fonds de seize
% mille feuillets.

\newcommand{\folder}[1]{\gdef\grothFolder{#1}}
\newcommand{\batch}[1]{\gdef\grothBatch{#1}}
\newcommand{\leaves}[2]{\gdef\grothFirst{#1}\gdef\grothLast{#2}}
\newcommand{\foldertitle}[1]{\gdef\grothFoldertitle{#1}}
\gdef\grothFolder{?}\gdef\grothBatch{?}\gdef\grothFirst{?}\gdef\grothLast{?}\gdef\grothFoldertitle{}

% --- Le résumé --------------------------------------------------------------
%
% Il ouvre la lecture modernisée, sous le titre « Résumé », et il est composé
% en petit. Ce n'est pas un ornement : c'est le seul passage du document qui
% ne soit pas des mathématiques, et un lecteur doit pouvoir le reconnaître
% avant de l'avoir lu — puis le sauter s'il connaît déjà le sujet, ou s'y
% arrêter s'il ne le connaît pas. Le filet qui le suit marque la reprise du
% corps.

\newenvironment{resume}
  {\small\setlength{\parskip}{0.45em}}
  {\par\medskip\hrule\medskip}

% --- Mots-clés ---------------------------------------------------------------
%
% En anglais, en clôture du résumé de la lecture modernisée : le vocabulaire
% moderne sous lequel chercher ce que les feuillets construisent. Le manifeste
% du site (`npm run manifest`) les extrait pour étiqueter la cote entière —
% cette ligne est la source unique des tags, il n'y a pas de fichier à côté.
\newcommand{\keywords}[1]{\par\smallskip\noindent
  {\footnotesize\textsc{Keywords} --- \textit{#1}}\par}

% --- L'appareil critique ----------------------------------------------------

% Le numéro de feuillet, en marge. C'est l'ancre qui permet de régler un
% désaccord sur une formule en regardant *un* feuillet plutôt qu'une chemise
% de six cents. Le site s'en sert aussi pour tourner les pages du fac-similé
% quand on fait défiler la transcription.
\newcommand{\leaf}[1]{%
  \marginnote{\footnotesize\color{gray!70}\textsf{#1}}%
  \label{leaf:#1}\ignorespaces}

% Illisible : on ne devine pas. Un mot inventé qui se lit comme les autres est
% le pire résultat possible.
\newcommand{\ill}{\textcolor{red!60!black}{[\,\ldots\,]}}

% Lecture douteuse : le mot est donné, et signalé comme incertain.
\newcommand{\uncertain}[1]{\uline{#1}}

% Ajout de l'éditeur — une lettre manquante, un mot sous-entendu.
\newcommand{\add}[1]{\textcolor{blue!50!black}{[#1]}}

% Biffé par Grothendieck lui-même. Ce qu'il a rayé fait partie du document :
% une rature dit souvent où la pensée a changé de direction.
\newcommand{\struck}[1]{\sout{#1}}

% Note du transcripteur, sur l'état matériel du feuillet.
\newcommand{\note}[1]{\par\smallskip{\small\color{gray!60!black}\textit{[#1]}}\par\smallskip}

% Note marginale de Grothendieck, distincte du corps du texte.
\newcommand{\marginal}[1]{\par\smallskip\noindent\hspace{1em}%
  {\small\color{gray!60!black}#1}\par\smallskip}

% --- Titre ------------------------------------------------------------------

\AtBeginDocument{%
  \begin{center}
    {\large\bfseries Fonds Alexandre Grothendieck --- cote n\textsuperscript{o}~\grothFolder}\\[2pt]
    {\itshape\grothFoldertitle}\\[4pt]
    {\small feuillets \grothFirst--\grothLast\ (lot \grothBatch)}\\[2pt]
    {\footnotesize\color{gray!60!black}Transcription établie d'après le fac-similé numérisé
     par l'Université de Montpellier.\\
     Les lectures incertaines sont soulignées, les illisibles notées [\,\ldots\,],
     les ajouts entre crochets.}
  \end{center}
  \vspace{1em}}

\endinput


\folder{137}
\batch{4}
\pages{61}{80}
\dating{[vers 1975-1976]}
\watermark{Édition de démonstration}
\foldertitle{Dévissage des complexes ½ simpliciaux ∂ - parfaits et structures multiplicatives… (1975 ou 1976) : notes manuscrites (s.d.)}

\begin{document}

\page{61}
\emph{Remarque.} Supposons $k$ \uncertain{principal}. Alors la catégorie
$\mathcal{P} \simeq \mathcal{P}$\,\ill{} \struck{\ill{}} admet des
\struck{$\Sigma$} noyaux (et des conoyaux, $\Sigma$ mais un
\uncertain{bimorphisme} n'est pas nécessairement un iso). Il
\struck{On peut \ill{} que les \ill{}}
\uncertain{est-il} vrai que si $N$ est le noyau de $A \to B$, alors
$N \to A$ est un $\Sigma$-mono, et s'insère donc dans une suite
$\Sigma$-exacte
\[
0 \to N \to A \to A' \to 0 ,
\]
mais le mono.\ $A' \to B$ n'est pas en général un $\Sigma$-mono.
\struck{Aussi} On ne peut \uncertain{interpréter}
$(\mathcal{P},$ \struck{\ill{}} $\mathbb{C}^{\bullet},$ \struck{\ill{}}$)$ comme
sol.\ d'un pb universel en \ill{} \ill{} les données précédentes :
\struck{Mais \ill{} \ill{} de façon naturelle, pour une $\mathcal{B}$
$k$-linéaire, la notion d'objet plat, et dont} le syst.\ ci-dessus est
universel pour les $\mathcal{B}$ \add{$k$-additives} munis d'un complexe
\struck{de cochaînes plats, \ill{} d'une \ill{} $\Sigma$ de
\ill{} exactes, telle que \ill{} tt $A \to B$ de $\mathcal{B}$ ait un
noyau tel que $\mathrm{Ker}\, u \to A$ soit \ill{} $\Sigma$)} $C^{\bullet}$
et d'une classe $\Sigma_C$ de suites exactes, telle que tt morphisme
$u : A \to B$ de $\mathcal{B}$ admette un noyau $K$ tel que $K \to A$ soit
un $\Sigma_C$-mono (on prend comme morphismes entre tels syst.\ ceux qui
transforment $C^{\bullet}$ en $C'^{\bullet}$ et qui sont
\emph{$\Sigma$-exacts} \ill{}).
\note{la page porte de longs passages biffés, encadrés et barrés de traits
ondulés ; on a donné ce qui s'en lit. « $k$-additives » est interlinéaire,
au-dessus de « munis ». Le premier « $\simeq$ » de la deuxième ligne est
d'une lecture douteuse.}

\page{62}
On travaillera désormais dans la catégorie $\mathcal{P}_1 = \mathcal{Q}$
(équivalente à la catégorie $\mathcal{Q}$ des complexes de chaînes finis de
$k$-modules \uncertain{tels} que les $L_i, Z_i, B_i, H_i$ sont libres de
type fini ; \uncertain{munie} de $\Sigma_1 = \Sigma \cap \mathcal{P}_1$),
dans laquelle la sous-catégorie des $\Sigma$-projectifs est
\uncertain{additive}, la sous-catégorie \add{additive} $\mathcal{P}_0$
engendrée par les $\mathbb{C}^i$ et $\Psi^0$ (i.e.\ par les $\Phi^i$), et celle des
$\Sigma$-injectifs est la sous-cat.\ additive $\mathbb{K}$
\struck{des \ill{}} \uncertain{engendrée} par les $\mathbb{C}^{\bullet}$ (les
$L_\bullet$ \struck{flasques} acycliques, i.e.\ les $L_\bullet$ flasques
\ldots).
\note{« additive » et « engendrée » sont écrits au-dessus de la ligne ; la
parenthèse « tels que \ldots\ $\Sigma_1 = \Sigma \cap \mathcal{P}_1$ » est
interlinéaire}
On s'intéressera aux \emph{résolutions droites $K^{\bullet}$
$\Sigma$-injectives de $\Psi^0$} (i.e.\ des résolutions flasques
$K^{\bullet}_{*}$ de $k_{*}$ \ill{} \uncertain{(qui donneront, dans un topos
algébrique,} des méthodes de calcul remarquables de la cohomologie à coeff.\
dans $k$). On a l'algèbre homologique « relative » qui donne des suites
exactes, qui nous disent que si (comme c'est le cas ici) tout objet
\struck{\ill{}} se plonge dans un $\Sigma$-injectif,
\struck{de façon} par un $\Sigma$-mono, alors \struck{les}
\struck{les deux résolutions \ill{} canoniquement \ill{}} deux résolutions
sont homotopes, et il y a une seule classe d'homotopie d'hom.\ de l'une
dans l'autre (\emph{NB} Il y a aussi assez de projectifs, et on a
l'énoncé analogue pour les résolutions projectives ; un \uncertain{fera}
attention que la cat.\ $\mathbb{K}$ des sommes de $\mathbb{C}^{\bullet}$
\uncertain{est}

\page{63}
contenue par \emph{assez} de $\Sigma$-projectifs ; il faut y rajouter
$\Psi^0 \ldots$) \struck{\ill{} \ill{}}
Les constructions sont d'ailleurs valables aussi dans la plus grande
catégorie $\mathcal{P}$ [définissant les $\Sigma$-proj.\ et les
$\Sigma$-inj., et le fait qu'il y en a assez.]
Les $\mathrm{Ext}^{i}$ calculés par les résolutions injectives ou
projectives \struck{sont} (à gauche \ldots) \uncertain{ou} les deux
simultanément, sont les $\mathrm{Ext}$ pour $i \geq 1$ ce sont les
hyperext habituels de $L_\bullet$ et de $M_\bullet$, (pour $i = 0$ ce sont
par contre les \uncertain{$Z^{i}$}.
\note{au-dessus de « injectives ou projectives », un mot interlinéaire,
peut-être « identiques », et « i.g. » au-dessus du « sont » biffé ; la
lecture de la phrase reste incertaine}

Si $F : \mathcal{Q} = \mathcal{P}_1 \subset \mathcal{P} \longrightarrow
\mathcal{B}$ ($k$-add.) est un foncteur, on peut construire les foncteurs
dérivés \add{droits ou gauches} (en utilisant des résolutions injectives ou
projectives). Nous sommes intéressés notamment aux
$\mathrm{R}^{i}F(\Psi^0)$, qu'on note aussi $\mathcal{H}^{i}(F)$ [on doit
pouvoir montrer que ce sont les dérivés \add{droits} de
$F \mapsto F(\Psi^0) = \Gamma(F)$ \add{sur les $\mathcal{P}$}, quand $F$
varie dans la cat.\ des foncteurs \add{($k$-add.)}
$\mathcal{P} \to \mathcal{B}$ \add{$\Sigma$}-exacts à gauche, (si on se
restreint à de tels $F$) ; que \uncertain{cette} catégorie est équivalente
à \struck{\ill{}} $\mathcal{B}^{\bullet}$, $F \mapsto F(\Psi^0)$ est
identifié au foncteur
\[
C^{\bullet} \longmapsto Z^{0}(C^{\bullet}) = H^{0}(C^{\bullet})
\quad \text{sur } \mathcal{B}^{\bullet},
\]
dont les dérivés sauf erreur sont les $H^{i}(C^{\bullet})$ pour
$i \geq 1$. On a aussi
\[
\mathrm{R}F(\Psi^0) \simeq F(\mathbb{C}^{\bullet})
\]
\note{les mots marqués comme ajouts sont interlinéaires, de sa main. Le
complexe universel est écrit tantôt avec un $C$ ordinaire, tantôt avec un
$C$ à double trait ; on l'a écrit $\mathbb{C}^{\bullet}$, comme aux pages
57 à 60, et $C^{\bullet}$ pour un complexe quelconque de $\mathcal{B}$.}

\page{64}
(iso.\ dans cat.\ dérivée) puisque $\mathbb{C}^{\bullet}$ est une résolution
$\Sigma$-injective de $\Psi^0$.
\note{un trait de séparation, au milieu de la page}

Nous considérons maintenant les \uncertain{réalisations} ½ simpliciales
\[
\mathcal{P}_{*} \supset \mathcal{Q}_{*} \supset \mathcal{P}_{0*} \supset
\mathbb{K}_{*}
\]
de $\mathcal{P} \supset \mathcal{Q} \supset \mathcal{P}_0 \supset
\mathbb{K}$. Or sur $\mathcal{P}_{*}$ il y a la structure $\otimes_{*}$ ou
$\overset{*}{\otimes}$ ($L_{*} \otimes_{*} M_{*}$ défini composante par
composante
\[
(L_{*} \otimes M_{*})_{[n]} = L_{[n]} \otimes M_{[n]} \ )
\]
avec contraintes d'ass., unité ($\Psi^0$), et de commutation,
\add{$\otimes$-bilin.} On note que $\mathbb{K}_{*}$ est un
\struck{\ill{} la réalisation (\ill{}) $\mathcal{P}^{*}, \mathcal{Q}^{*}$}
\uncertain{« idéal »} (i.e.\ si un des facteurs $L_{*}, M_{*}$ est dans
$\mathbb{K}_{*}$, \uncertain{l'autre} aussi), $\mathcal{P}_0$ et
$\mathcal{Q}$ sont stables par $\otimes$, car
\[
\underset{\mathbf{1}_{*}}{\Psi^{0}_{*}} \otimes
\underset{\mathbf{1}_{*}}{\Psi^{0}_{*}} \simeq
\underset{\mathbf{1}_{*}}{\Psi^{0}_{*}}
\quad (\Rightarrow \text{stabilité de } \mathcal{P}_0)
\]
et
\[
\Psi^{i}_{*} \otimes \Psi^{j}_{*} \simeq \Psi^{i+j}_{*} +
\text{élément de } \mathbb{K}
\quad (\Rightarrow \text{stabilité de } \mathcal{P}_{1*} = \mathcal{Q})
\]
\marginal{\ill{} $\mathcal{P}_0$ \ill{} \ill{} (\ill{})}
\note{sous chacun des trois $\Psi^{0}_{*}$, un signe $=$ vertical et
$\mathbf{1}_{*}$. La note marginale, écrite de biais dans la marge gauche
en face des lignes biffées, est en partie barrée et presque entièrement
illisible.}

On transporte ces structures multiplicatives sur $\mathcal{P}$ et ses
sous-catégories (notation $*$, pour éviter des confusions) et ses autres
réalisations $\mathcal{P}_{\bullet}$ et $\mathcal{P}^{*}$,
$\mathcal{P}^{\bullet}$ (au sens des \ill{} \uncertain{contravariantes}).
On note $*$ sur $\mathcal{P}_{\bullet}$ et $\mathcal{P}^{\bullet}$ et
$\overset{*}{\otimes}$ sur $\mathcal{P}^{*}$, car l'équivalence \ill{} de

\page{65}
$\mathcal{P}_{*}$ avec $\mathcal{P}^{*}$ est compatible avec $\otimes$.

Soit $\mathcal{B}$ une cat.\ ($k$-)additive avec op.\ ($k$-)bil.\
$\otimes$ \add{(propriété dans chaque facteur $\otimes$)} munie d'une
contrainte \add{(données)} d'ass., d'unité, et de comm.\ \struck{\ill{}}
satisfaisant aux compatibilités habituelles (propriétés), soit
$\mathcal{B}'$ une autre telle. Si on a un foncteur ($k$-)additif
(propriété)
\[
F : \mathcal{B}' \longrightarrow \mathcal{B}
\]
\uncertain{une} structure \emph{multiplicative} sur $F$ \ill{} homs de
foncteurs
\[
F(X) \otimes F(Y) \xrightarrow{\ u_{X,Y}\ } F(X \otimes Y)
\]
et iso \struck{\ill{}}
$\mathbf{1}_{\mathcal{B}} \xrightarrow{\ \alpha_F\ } F(\mathbf{1}_{\mathcal{B}'})$
qui soient compatibles aux données d'ass., d'unité et comm. :

\begin{tikzcd}[column sep=small, row sep=normal, nodes={font=\small}]
(F(X) \otimes F(Y)) \otimes F(Z) \arrow[r, "{u_{X,Y} \otimes F(Z)}"] \arrow[d, "\wr\ a^{\otimes}_{F(X),F(Y),F(Z)}"'] & F(X \otimes Y) \otimes F(Z) \arrow[d, "u_{X \otimes Y,Z}"] \\
F(X) \otimes (F(Y) \otimes F(Z)) \arrow[d, "F(X) \otimes u_{Y,Z}"'] & F((X \otimes Y) \otimes Z) \arrow[d, "\wr\ F(a^{\mathcal{B}'}_{X,Y,Z})"] \\
F(X) \otimes F(Y \otimes Z) \arrow[r, "u_{X,Y \otimes Z}"'] & F(X \otimes (Y \otimes Z))
\end{tikzcd}

\begin{tikzcd}[column sep=normal, row sep=normal, nodes={font=\small}]
F(X) \otimes \mathbf{1}_{\mathcal{B}} \arrow[r, "F(X) \otimes \alpha"] \arrow[ddr, "\wr\ \delta^{\mathcal{B}}_{F(X)}"'] & F(X) \otimes F(\mathbf{1}_{\mathcal{B}'}) \arrow[d, "u_{X,\mathbf{1}_{\mathcal{B}'}}"] \\
 & F(X \otimes \mathbf{1}_{\mathcal{B}'}) \arrow[d, "\wr\ F(\delta^{\mathcal{B}'}_{X})"] \\
 & F(X)
\end{tikzcd}

et symétrique

\begin{tikzcd}[column sep=large, row sep=normal, nodes={font=\small}]
F(X) \otimes F(Y) \arrow[r, "u_{X,Y}"] \arrow[d, "\sigma^{\mathcal{B}}_{F(X),F(Y)}"'] & F(X \otimes Y) \arrow[d, "F(\sigma^{\mathcal{B}'}_{X,Y})"] \\
F(Y) \otimes F(X) \arrow[r, "u_{Y,X}"'] & F(Y \otimes X)
\end{tikzcd}

\note{dans chacun des trois diagrammes, l'auteur écrit « comm. » au centre.
Au premier, le but de la flèche horizontale supérieure est corrigé sur la
page.}

\page{66}
\emph{NB} $F(\mathbf{1}_{\mathcal{B}'})$ devient donc un « objet-\struck{\ill{}}anneau »
de $\mathcal{B}$ (« $k$-Algèbre » quand il y a un $k$).
$F$ transforme « anneaux » de $\mathcal{B}'$ en « anneaux » de $\mathcal{B}$,
en respectant les propriétés d'unité, associativité et commutativité.

Les $F$ en question forment une catégorie. Si on a $\mathcal{B}''$,
$\mathcal{B}'$, $\mathcal{B}$, ils se composent, d'où
2-catégorie des $\otimes$-catégories $k$-additives.

\emph{Supposons maintenant que $\mathcal{B}' = \mathbb{K}$}
(catégorie des coquilles de \uncertain{Chaînes} \uncertain{des} résolutions
injectives \ldots). Un foncteur $k$-lin.
\[
\mathbb{K} \xrightarrow{\ F\ } \mathcal{B}
\]
s'identifie donc à un $\mathbb{C} \in \mathcal{B}^{\bullet}$.
Donc un sens à ce qu'il faudrait entendre par « structure multiplicative »
sur le complexe de cochaînes $\mathbb{C}^{\bullet}$ de $\mathcal{B}$
(on dira : str.\ mult.\ « demi-simpliciale » pour ne pas prêter à
confusion).
\note{« coquilles » : lecture d'après la page ; la fin de la parenthèse
est illisible en partie}

\page{67}
comme \emph{donnée}, c'est-à-dire \uncertain{la} donnée des
hom.\ fonctoriels
\[
F(X) \otimes F(Y) \longrightarrow F(X \otimes Y)
\]
où il suffit de se donner sur la \struck{\ill{}} catégorie
\add{$k$-lin.} $\mathbb{K} \times \mathbb{K}$ \add{de $\mathbb{K} \times
\mathbb{K}$} (dans $\mathbb{C} \times \mathbb{C}$,
avec sa structure $k$-lin.)
\struck{\ill{}}, donc de se donner des homs de bicomplexes de cochaînes
\[
F(\mathbb{C}^{\bullet}) \otimes F(\mathbb{C}^{\bullet}) \longrightarrow
F(\mathbb{C}^{\bullet} * \mathbb{C}^{\bullet}) \quad \text{dans }
\mathcal{B}^{\bullet\bullet}
\]
\note{sous les deux premiers $F(\mathbb{C}^{\bullet})$, un signe $=$
vertical, $\mathbb{C}^{\bullet}$, et un $\mathcal{B}$ entre eux, sous le
$\otimes$}
ce qui \uncertain{peut} \uncertain{s'expliciter}, si
$\mathbb{C}^{\bullet}$ \struck{$\otimes$} $\mathbb{C}^{\bullet}$ correspond au
\add{bi}complexe de chaînes \add{défini par}
\[
\mathbb{C}^{i}_{\bullet} * \mathbb{C}^{j}_{\bullet} =
\mathrm{N}_{\bullet}(\mathbb{C}^{i}_{*} \otimes
\mathbb{C}^{j}_{*})
\]
\note{après le signe $=$, un symbole biffé, illisible ; sous les deux derniers facteurs, $= \overset{i}{\Lambda}\Phi_{*}
\otimes \overset{j}{\Lambda}\Phi_{*}$ (lecture incertaine des exposants)}
\[
\mathbb{F}^{ij} \qquad F(\mathbb{C}^{i}) \otimes_{\mathcal{B}}
F(\mathbb{C}^{j}) \longrightarrow
\underline{\mathrm{Hom}}_{k}\bigl((\mathbb{C}^{i}_{\bullet} *
\mathbb{C}^{j}_{\bullet})^{\vee}, C^{\bullet}\bigr)
\]
\note{sous le membre de gauche, $= C^{i}$ et $= C^{j}$ ; sous
$(\mathbb{C}^{i}_{\bullet} * \mathbb{C}^{j}_{\bullet})$,
« $= C_i * C_j$ » (lecture incertaine)}
\marginal{\ill{} \ill{} \ill{}}
\note{dans la marge gauche, une note de biais de quatre ou cinq mots,
illisible}
i.e.
\[
\mathbb{F}^{ij} \qquad (C^{i} \underset{\mathcal{B}}{\otimes} C^{j})
\underset{k}{\otimes}
\underbrace{(C^{\bullet}_{i} * C^{\bullet}_{j})}_{\in\, \mathbb{K}^{\bullet}}
\longrightarrow C^{\bullet}
\]
\note{au-dessus de $(C^{\bullet}_{i} * C^{\bullet}_{j})$, une accolade
avec « $C^{\bullet}_{ij}$ » et un mot abrégé, peut-être « défini » ; le
préfixe écrit devant les deux formules, lu $\mathbb{F}^{ij}$, est d'une
lecture douteuse}
avec compatibilité avec \uncertain{les} diff.\ \ill{} \ill{}
\uncertain{variable}.

On peut aussi l'écrire \struck{$(C^{i} \otimes_{\mathcal{B}} C^{j}) \otimes_{k}$}
\begin{align*}
(1^{*}) \qquad & C^{i} \underset{\mathcal{B}}{\otimes} C^{j}
\longrightarrow \underline{\mathrm{Hom}}^{*}_{k}(C^{*}_{i} \otimes
C^{*}_{j},\ C^{*}) \\
\text{ou encore} \qquad
(2^{*}) \qquad & (C^{i} \otimes_{\mathcal{B}} C^{j}) \otimes_{k}
\underbrace{(\overset{i}{\Lambda}\Phi^{*} \otimes
\overset{j}{\Lambda}\Phi^{*})}_{C^{*}_{ij}}
\xrightarrow{\ u^{ij}\ } C^{*}
\end{align*}
\note{sous $C^{*}_{i} \otimes C^{*}_{j}$ dans $(1^{*})$ :
$= \overset{i}{\Lambda}\Phi^{*} \otimes \overset{j}{\Lambda}\Phi^{*}$ ;
sous $C^{*}$ : $= \mathrm{DP}^{*}(C^{\bullet})$. Un double trait vertical
dans la marge gauche, en face de $(1^{*})$ et $(2^{*})$.}

\page{68}
homs compatibles avec op.\ ½ simpliciales ($i, j$ fixés) et avec op.\ diff.\
(pour $i, j$ variables)

\begin{tikzcd}[column sep=large, row sep=normal, nodes={font=\small}]
(C^{i} \otimes_{\mathcal{B}} C^{j}) \otimes C^{*}_{ij} \arrow[d] & (C^{i} \otimes_{\mathcal{B}} C^{j}) \otimes C^{*}_{i+1,j} \arrow[l, "d^{*}_{i+1,j}"'] \arrow[d, "d^{i} \otimes -\,\otimes -"] \\
C^{*} & (C^{i+1} \otimes_{\mathcal{B}} C^{j}) \otimes C^{*}_{i+1,j} \arrow[l, "d"]
\end{tikzcd}

\note{sur la page, la flèche verticale de droite n'est indiquée que par son
étiquette ; au début de la ligne inférieure, une écriture raturée et
illisible précède $C^{*}$}
plus des données de compatibilités avec unité, ass., comm.
(\uncertain{Elles} devraient s'expliciter de façon particulièrement triviale).

\emph{Remarque} $\mathbb{K}$ n'a pas d'unité pour $*$\,! donc on ne peut
exprimer directement, en restant dans $\mathbb{K}$, une compatibilité avec
les \struck{unités}\,! Comment s'en tirer\,? Si on passe à $\mathcal{P}_0$
ou $\mathcal{Q} = \mathcal{P}_1$,
\struck{\ill{} $\mathcal{B}$ \ill{} \ill{} \ill{} $\Sigma$ \ill{}}
on trouve que les foncteurs $k$-additifs de $\mathcal{Q}$ (resp.\
$\mathcal{P}_0$) dans $\mathcal{B}$ qui sont $\Sigma$-exacts à gauche
(i.e.\ transforment \struck{des} $\Sigma$-suites exactes courtes
$0 \to A' \to A \to A'' \to 0$ en suites exactes
$0 \to F(A') \to F(A) \to F(A'')$\,) \struck{\ill{}}
\struck{$F(A') \to F(A) \ldots \Sigma$ \ill{}} correspondent par
$F \mapsto F(\mathbb{C}^{\bullet})$ aux complexes de cochaînes dans
$\mathcal{B}$ tels que $\forall i$ \uncertain{Coker}
\note{un double trait vertical dans la marge droite, en face des dernières
lignes}

\page{69}
($i \geq 0$ seulement) $C^{i} \xrightarrow{\ d^{i}\ } C^{i+1}$ ait un
noyau \struck{qui soit tel que $Z^{i} \to C^{i}$ soit un $\Sigma$-mono
($\mathcal{B}$ \ill{} de $\Sigma$ dans $\mathcal{B}$ \ill{}) de la
catégorie ambiante, il n'y a \ill{} \ill{} $\ldots$}
(ceci dit, on a bien
\note{ce passage est encadré et barré de plusieurs traits obliques}
à de tels complexes, qui correspondent \add{($\Sigma$-ex.\ à g.)}
\struck{\ill{}} \uncertain{$k$-add.}
\struck{donc à des foncteurs \ill{}
$\mathcal{P}_0 \to \mathcal{B}$ resp.\ $\mathcal{Q} \to \mathcal{B}$.
On a \ill{} \add{\ill{}} donner une structure multiplicative sur un
\ill{}, ce qui nous oblige à donner de plus
$F(\Psi^0) \otimes_{\mathcal{B}} F(\mathbb{C}^{i}) \to
F(\Psi^0 * \mathbb{C}^{i})$ et \ill{}}
\struck{\ill{} $\mathcal{B}$ bifonctoriels
$F(X) \otimes_{\mathcal{B}} F(Y) \longrightarrow F(X * Y)$
sur $\mathcal{P}_0 \times \mathcal{P}_0$ resp.\
$\mathcal{Q} \times \mathcal{Q}$. On}
\note{une longue accolade dans la marge gauche embrasse ces lignes, que
barrent aussi trois longs traits obliques ; en face, dans la marge, « $*$ »
et un point d'exclamation}
\add{\uncertain{Or}} si \uncertain{tout} dans $\mathcal{B}$ \uncertain{a}
\uncertain{des} \uncertain{facteurs} \struck{\ill{}} \ill{}
\uncertain{conoyaux}, alors on a \uncertain{trouvé} \uncertain{les}
foncteurs $\Sigma$-ex.\ à g.\ de $\mathcal{P}$ dans $\mathcal{B}$
\struck{\ill{}} \uncertain{interviennent} aussi comme \ill{}
\uncertain{générales} des foncteurs $k$-additifs
\uncertain{associés} à $C^{\bullet} \in \mathcal{B}^{\bullet}$. On peut
\uncertain{les} \uncertain{ré}\ill{} \uncertain{explicitement} ainsi :
$\forall\, L \in \mathcal{P}$ \struck{\ill{}}
\add{($\in \mathcal{Q}$, $\in \mathcal{P}_0$)} s'insère dans deux
$\Sigma$-\uncertain{suites} de $\mathcal{P}$ ($\mathcal{Q}$,
$\mathcal{P}_0$)
\[
\begin{array}{l}
0 \to L \to I' \to L'' \to 0 \\
\qquad\qquad 0 \to L' \to J'' \to L'' \to 0
\end{array}
\]
avec $I, J \in \mathbb{K}$, et par suite on \ill{}
\note{la page est d'une écriture rapide, et le texte non biffé est d'une
lecture très incertaine ; le second terme de la deuxième suite est
surchargé}

\page{70}
Calcul \uncertain{$F$} : \uncertain{l'avoir} \ill{} \ill{} \ill{}
$\mathbb{K}$ comme
\[
F(X) \simeq \mathrm{R}^{0}F(L) \simeq \mathrm{Ker}\bigl(F(I) \to F(J)\bigr)
\]
Ceci dit, \ill{} on a

\begin{tikzcd}[column sep=large, row sep=small]
F(X) \otimes_{\mathcal{B}} F(Y) \arrow[r, "?"] \arrow[d] & F(X * Y) \arrow[d] \\
F(I) \otimes_{\mathcal{B}} F(J) \arrow[r] & F(I * J)
\end{tikzcd}

\note{sous la flèche horizontale inférieure : « comme \ill{} par $F \mid
\mathbb{K}$ et la structure multi. »}
Or on sait que \struck{si} $*$ transforme des couples de $\Sigma$-monos en
un $\Sigma$-mono (résulte immédiatement du \uncertain{théorème} de
Dold-Puppe) donc la 2\textsuperscript{e} flèche verticale est un mono,
\uncertain{unique} donc\,? \ill{} la dite \uncertain{factorisation}
\ill{} \ill{} (\ill{} fait qu'il \uncertain{soit} \uncertain{naturel}
\uncertain{et} \uncertain{aussi} deux faits \ill{} \uncertain{immédiats},
\uncertain{et} aussi qu'il satisfait aux compatibilités d'ass.\ et de
comm.\ si la restriction à $\mathbb{K}$ y satisfait. Il reste donc
l'ex.\ unité ; \ill{} donné
\struck{qu'il suffit que}
\[
\underline{\mathbf{1}}_{\mathcal{B}'} \longrightarrow
F(\Psi_0) \simeq Z^{0}(C^{\bullet})
\]
\note{dans la formule, un $F(\Psi_0)$ biffé après $\underline{\mathbf{1}}_{\mathcal{B}'}$, et un symbole biffé, illisible, avant $Z^{0}$}
et on sait que $C^{\bullet}$ \uncertain{a} \uncertain{une}
\emph{\uncertain{augmentation}} \ill{} \ill{} \ill{} unité \ill{} que
\ill{} \ill{} \ill{} « $k$-$\otimes$-alg » de $F(\Psi^0)$ ; la structure
de « $k$-$\otimes$-alg.\ commutative » (sur \ill{} \ill{} \ill{}
$\mathbb{K}$)
\note{les dernières lignes, au bas du feuillet, sont d'une écriture
serrée et en partie illisible}
\marginal{\emph{NB} On \ill{} \ill{} que \ill{} besoin de l'existence de
$\mathrm{Ker}(d^0 : C^0 \to C^1)$, donc \ill{} \ill{} « il suffit de dire
que $\mathbf{1}_{\mathcal{B}} \to C^0$ \ill{} \ill{} $C^0 \to C^1$ \ill{}
\ill{} \ill{} \ill{} et qui \ill{} \ill{} \ill{} \ill{} \ill{}
\struck{\ill{}} $C^{0}$ \ill{} \ill{} »}
\note{cette note, écrite verticalement dans la marge gauche sur toute la
hauteur de la page, est reliée par un trait à la dernière ligne ; la
lecture en est très fragmentaire}

\page{71}
\emph{Applications} Toutes les constructions (« universelles ») dans
$\mathcal{P}_{*}$ \struck{qui} (ou $\mathcal{Q}_{*}$, $\mathcal{P}_{0*}$,
$\mathbb{K}_{*}$) qui ne font intervenir que la structure $k$-lin.\ et
\add{les} \uncertain{multiples} \add{\ill{}} $\otimes$-structures,
\struck{donnent} \uncertain{donnent} des constructions analogues sur un
$C^{\bullet}$ ($\in \mathcal{B}^{\bullet}$ muni de structure
multiplicative).
\note{un double trait vertical dans la marge gauche, en face des premières
lignes}

\emph{Ex.} \struck{$\Psi^{0}$ \ill{} \ill{} \ill{}} Comme
$\mathbb{C}^{\bullet}_{*} \otimes \mathbb{C}^{\bullet}_{*}$ est une
résolution $\Sigma_{*}$-inj.\ de $\Psi^{0}_{*}$ dans $\mathcal{P}_{0*}$,
et comme $\mathbb{C}^{\bullet}_{*}$ l'est aussi, on trouve
\[
\mathrm{Tot}^{\bullet}(\mathbb{C}^{\bullet}_{*} * \mathbb{C}^{\bullet}_{*})
\longrightarrow \mathbb{C}^{\bullet}_{*}
\]
\marginal{\uncertain{défini à homotopie}}
près, d'où un hom
\[
\begin{array}{ccc}
F_{C^{\bullet}}(\mathbb{C}^{\bullet}_{*} * \mathbb{C}^{\bullet}_{*}) &
\longrightarrow & F(\mathbb{C}^{\bullet}_{*}) \\
\uparrow & & \Vert \\
F_{C^{\bullet}}(\mathbb{C}^{\bullet}) \otimes_{\mathcal{B}}
F_{C^{\bullet}}(\mathbb{C}^{\bullet}) & & C^{\bullet} \\
\Vert & & \\
C^{\bullet} \otimes C^{\bullet} & &
\end{array}
\]
défini à homotopie près, d'où en composant un hom défini à homotopie près
\[
C^{\bullet} \otimes_{\mathcal{B}} C^{\bullet} \longrightarrow C^{\bullet}
\]
avec conditions d'associativité, de commutation (attention aux signes\,!)
et d'unité --- le tout à homotopie près. Il en résulte une structure de
« $\otimes$-$\mathcal{B}$-alg.\ » \uncertain{graduée} sur
$H^{*}(C^{\bullet})$, avec les propriétés habituelles.
\note{la note marginale, écrite sous la flèche
$\mathrm{Tot}^{\bullet}(\ldots) \to \mathbb{C}^{\bullet}_{*}$, y est
reliée par une flèche}

\page{72}
\emph{Ex 2} Le « complexe de De Rham à p.d.\ » est défini \struck{comme}
\uncertain{comme} \uncertain{objet} \ill{} (\uncertain{avec} \ill{} les
p.d.) \struck{comme un objet} objet bigradué de
\struck{\ill{}} $\mathcal{Q}_{*}$, avec structure d'algèbre
\emph{bigraduée} \uncertain{standard} \ill{}. Donc $F_{C^{\bullet}}$ le
transforme en objet de \ill{} \ill{}. On ne peut \uncertain{affirmer}
encore (sans introduire la « structure p.d. » sur $F$) que le complexe de
De Rham associé à $F$ est \uncertain{stable} (\ill{} \ill{} \ill{},
\uncertain{etc.}). Le \struck{fait} que
\struck{\ill{} \ill{} \ill{}} \ill{} \ill{} \ill{} \ill{}\,), on trouve
une résolution de $k_{*} = \Psi^{0}_{*}$, flasques (\uncertain{resp.}\
$\Sigma$-injectives) sauf en degré \uncertain{maximum} $n$ (où on trouve
$\Psi^{n}_{*} = \overset{n}{\Lambda}\Psi_{*}$) --- donc qui \ill{} une
résolution $\Sigma$-injective tronquée, \uncertain{montre} que le
\ill{} de sa transformée par $F$ \uncertain{proprement} \uncertain{dite},
et \uncertain{coïncide} avec celle de $C^{\bullet}$ \struck{\ill{}}
\ill{} \ill{} \ill{} \ill{}, le complexe \uncertain{est} homotopes (dans
\uncertain{classe} d'homotopie \uncertain{donnée}) au complexe tronqué
\uncertain{naïf} de $C^{\bullet}$ --- en \uncertain{fait} \ill{} un choix
\uncertain{canonique} (car $\mathrm{DR}^{n}_{*}
\xrightarrow{\ \deg \leq n\ } \mathbb{C}^{\bullet}_{*}$ hom.\ unique
\emph{\uncertain{\ill{}}}).
\note{un double trait vertical dans la marge droite, en face des premières
lignes. La page est d'une écriture rapide, et la syntaxe de la longue
phrase qui suit « Le fait que » ne se laisse pas établir ; les mots
biffés au début de la neuvième ligne sont soulignés.}

\page{73}
Il y a donc un complexe de De Rham-Sullivan \struck{\ill{}}
\uncertain{si} \uncertain{$k$} \uncertain{de car.}\ $0$.
\note{« si $k$ de » est écrit au-dessus du mot biffé ; le reste du
feuillet est blanc}

\page{74}
\subsection*{Catégories polynomiales}
\note{le titre est souligné sur la page ; il ouvre une nouvelle suite, dont
les axiomes a), b), c), \ldots\ sont ceux que cite la page 83 (« les axiomes
a) à f) »)}

Une \emph{catégorie polynomiale} est une catégorie $\mathcal{P}$, munie
d'une sous-catégorie (\uncertain{non} pleine en général)
$\mathcal{P}_1$, avec les axiomes suivants
\marginal{\uncertain{notation} $\mathcal{P}(X, Y)$ \uncertain{pour}
$\mathrm{Hom}_{\mathcal{P}}(X, Y)$}

a) $\mathcal{P}$ admet des produits finis, \struck{et l'objet final
\ill{}}

b) $\mathrm{Ob}\,\mathcal{P}_1 = \mathrm{Ob}\,\mathcal{P}$,
\struck{les produits et pour deux objets dans $\mathcal{P}_1$, leur
produit dans $\mathcal{P}$ est un produit dans $\mathcal{P}_1$ (i.e.\ si
$X, Y, Z \in \mathrm{Ob}\,\mathcal{P}_1$, \ill{} \ill{} \ill{}
$Z \to X \times Y$ de $\mathcal{P}$ est dans $\mathcal{P}_1$ ssi
$u = (v, w) : X \to Y \times Z$ est dans $\mathcal{P}_1$ \ill{}}
\note{ce passage, surchargé d'ajouts interlinéaires (« i.e.\ si
$X, Y, Z \in \mathrm{Ob}\,\mathcal{P}_1$ », « $\mathcal{P}$-flèches »), est
encadré et barré de traits obliques}
(donc $\mathcal{P}_1$ est décrite par un \struck{\ill{}} ensemble
$\mathrm{Fl}_1 \subset \mathrm{Fl}(\mathcal{P}) = \mathrm{Fl}$ contenant
les identités, et stable par composition) On veut de plus que
\struck{$\mathrm{Fl}_1$ contienne les iso, et que} pour
$X, Y \in \mathrm{Ob}\,\mathcal{P}$, leur produit $X \times Y$ soit un
produit dans $\mathcal{P}_1$ (i.e.\ une flèche
$u = (v, w) : Z \to X \times Y$ est « de degré $1$ » ssi ses composantes
$v$, $w$ le sont), enfin \uncertain{que} $0$ \uncertain{est} objet final
\add{de $\mathcal{P}$} et $X \to 0$ est de degré~$1$.

c) $\mathcal{P}_1$ est une catégorie \add{(donc) additive} \ill{}
\uncertain{de} $\mathcal{P}_1$.)

Cela implique que \struck{\ill{} objet \ill{} \ill{} \ill{}
\ill{}} \add{les $\mathrm{Hom}_{\mathcal{P}}(X, Y)$}
\struck{$\mathbb{Z} \to \mathrm{End}(\mathrm{id}_{\mathcal{P}_1}) \to
\mathrm{End}(\mathrm{id}_{\mathcal{P}})$}
\struck{\ill{} \ill{} \ill{} $X, Y \in \mathrm{Ob}\,\mathcal{P}$, \ill{}
\ill{} $X$} sont des groupes abéliens (car les objets $X$ \add{de
$\mathcal{P}$} sont \struck{\ill{}} des objets groupes abéliens de
$\mathcal{P}$), dépendant fonctoriellement de $X$, mais pas de $Y$
(i.e.\ on a
\[
\boxed{(f + g) \circ h = f \circ h + g \circ h}
\]
\note{« i.e.\ pour $X$ » est écrit sous la ligne biffée ; le passage
biffé, où figure la ligne
$\mathbb{Z} \to \mathrm{End}(\mathrm{id}\,\mathcal{P}_1) \leftarrow
\mathrm{End}(\mathrm{id}\,\mathcal{P})$, est entouré d'une grande courbe
qui le relie au reste de la page}
mais \struck{\ill{}} \struck{\ill{}} \uncertain{on n'aura pas}
éventuellement
\[
f \circ (g + h) \neq f \circ g + f \circ h \ !
\]
Il en résulte aussi que $\mathbb{Z}^{\times}$ opère sur $\mathcal{P}$
\uncertain{i.e.\ on a} un \struck{\ill{}} \uncertain{hom} multiplicatif
\[
\mathbb{Z}^{\times} \longrightarrow \mathrm{End}(\mathcal{P}_1)
\longleftarrow \mathrm{End}(\mathcal{P})
\]
pour $n \in \mathbb{Z}$ et $X \in \mathrm{Ob}\,\mathcal{P}_1 =
\mathrm{Ob}\,\mathcal{P}$ \uncertain{associe} $n\,\mathrm{id}_X$.
Ainsi $\mathbb{Z}$ opère dans \struck{\ill{}}
$\mathcal{P}(X, Y)$ par $\mu_n(f) = f \circ (n\,\mathrm{id}_Y)$
\note{sic : $\mathrm{id}_Y$ sur la page, où l'on attendrait
$\mathrm{id}_X$ ; les deux $\mathrm{End}$ portent sur
$\mathrm{id}_{\mathcal{P}_1}$, $\mathrm{id}_{\mathcal{P}}$ dans la ligne
biffée plus haut}

\page{75}
d) Il existe une décomposition de
\struck{$\mathrm{Hom}$}\,$\mathcal{P}(X, Y)$, en
\[
\mathcal{P}(X, Y) = \bigoplus_{n \geq 0} \mathcal{P}_n(X, Y)
\]
telle que pour $f_n \in \mathcal{P}_n(X, Y)$, et $\lambda \in \mathbb{Z}$,
on ait
\[
f_n \circ (\lambda\,\mathrm{id}_X) = \lambda^{n} f_n
\]
Plus généralement, si $X_1, \ldots, X_r, Y$ sont des objets de
$\mathcal{P}$, on a une décomposition
\[
(*) \qquad \mathcal{P}(X_1 \times \cdots \times X_r, Y) =
\bigoplus_{n = (n_1, \ldots, n_r) \in \mathbb{N}^{r}}
\mathcal{P}_n(X_1, \ldots, X_r; Y)
\]
avec, pour $f \in \mathcal{P}_{n_1, \ldots, n_r}(X_1, \ldots, X_r; Y)$ et
$\lambda = (\lambda_i)_{1 \leq i \leq r} \in \mathbb{Z}^{r}$, on ait
\[
f \circ (\lambda_1 \mathrm{id}_{X_1}, \ldots, \lambda_r \mathrm{id}_{X_r})
= \lambda_1^{n_1} \cdots \lambda_r^{n_r} f
\]
\note{un double trait vertical dans la marge gauche, en face de d). Le
$\mathcal{P}$ de $(*)$ est récrit sur un premier signe.}

\emph{Remarque} Cette décomposition \uncertain{est} \uncertain{unique},
\uncertain{en} \uncertain{vertu} de \uncertain{considérations}
« arithmétiques » --- il le faudrait\,!

\emph{Lemme} Soit $P$ un
groupe abélien, et $u : (\mathbb{Z}^{\times})^{r} \to \mathrm{End}(P)$ un
hom de monoïdes. \struck{$P \ldots \mathbb{A}$ \ldots\
$n = (n_1, \ldots, n_r) \in \mathbb{Z}$} Supposons qu'il existe une
décomposition en somme directe
\[
P = \sum_{n = (n_1, \ldots, n_r) \in \mathbb{N}^{r}} P_n
\]
telles que pour $f \in P_n$, et $\lambda = (\lambda_1, \ldots, \lambda_r)
\in \mathbb{Z}^{r}$, on ait
\[
u(\lambda) \cdot f = \lambda_1^{n_1} \cdots \lambda_r^{n_r} f
\]
Alors cette décomposition est unique, plus précisément on a pour tout
$n \in \mathbb{N}^{r}$
\[
P_n = \lbrace f \in P \mid u(\lambda) f = \lambda_1^{n_1} \cdots
\lambda_r^{n_r} f \quad \forall \lambda \in \mathbb{Z}^{r} \rbrace
\]
Il suffit de voir que si $n' \neq n$ et si $g_{n'} \in P_{n'}$ est tel que
$u(\lambda) g_{n'} = \lambda_1^{n_1} \cdots \lambda_r^{n_r} g_{n'}$
$\forall \lambda \in \mathbb{Z}^{r}$, alors $g_{n'} = 0$. Or on a aussi
$u(\lambda) g_{n'} = \lambda_1^{n'_1} \cdots \lambda_r^{n'_r} g_{n'}$
\note{du « Lemme » jusqu'ici, tout est barré de quatre longs traits
obliques ; on le donne tel qu'il se lit, sans le marquer comme biffé}
\marginal{\emph{Ex} Si $P$ est un groupe annulé par un nombre premier
$p$, \ill{} $\lambda \mapsto \lambda^{p}\,\mathrm{id}_{\mathbb{Z}}$ et
$\lambda \mapsto \lambda\,\mathrm{id}$ \uncertain{coïncident} ; \ill{}
\ill{} \ill{} $1$, et $p$\,!}
\note{la note, écrite verticalement dans la marge gauche, est donnée ici ;
à côté, dans la même marge, une colonne de signes épars (accolades,
$\Sigma$, $\mathbb{N}$) sans suite lisible}

\page{76}
donc
\[
(\lambda^{n} - \lambda^{n'})\, g_{n'} = 0 \qquad \forall \lambda \in
\mathbb{Z}^{r}
\]
et il suffit de voir que les $\lambda^{n} - \lambda^{n'}$ engendrent
l'idéal $1$, \uncertain{ce} \uncertain{qui} \ill{}
\note{ces premières lignes, qui achèvent le lemme de la page 75, sont
barrées comme lui}
\uncertain{Si} on exigeait que, pour tout $\mathbb{Z}$-algèbre $k$, les
\uncertain{composantes} \uncertain{entre} les $\mathcal{P}_k(X, Y) =
\mathcal{P}(X, Y) \otimes_{\mathbb{Z}} k$, et on exigeait la
décomposition plus haut valable \struck{pour} dans tt
$\mathcal{P}_k(\prod X_i, Y)$, pour $\lambda \in (k^{\times})^{r}$. On va
\uncertain{mettre} comme \emph{structure} supplémentaire la \emph{donnée}
de la décomposition $(*)$, avec

e) Condition de degré pour la composition de
\[
\begin{array}{ccccc}
X_1 & \times \cdots \times & X_r & \xrightarrow{\ (n_1, \ldots, n_r)\ } &
Y \\
\uparrow {\scriptstyle m^{1}_{1}, \ldots, m^{1}_{s_1}} & &
\uparrow {\scriptstyle m^{r}_{1}, \ldots, m^{r}_{s_r}} & & \\
Z^{1}_{1} \times \cdots \times Z^{1}_{s_1} & & Z^{r}_{1}, \ldots,
Z^{r}_{s_r} & &
\end{array}
\]
le composé est de degré $(n_i m^{i}_{1}, n_i m^{i}_{2}, \ldots,
n_i m^{i}_{s_i})$ sur $Z^{i}_{1}, \ldots, Z^{i}_{s_i}$ ; donc

f) condition \add{\uncertain{de symétrie}} \struck{\ill{}} de permutation
des facteurs

g) $\forall\, X \in \mathrm{Ob}\,\mathcal{P}$, l'application canonique
déduite de $X \to 0$
\[
\mathcal{P}_0(0, Y) \longrightarrow \mathcal{P}_0(X, Y)
\]
est \uncertain{bijective} (les « applications constantes » de $X$ dans
$Y$ « ne dépendent pas de $X$ »)
\marginal{\uncertain{Oublié} : \ill{} \ill{} \ill{} \ill{} \ill{}
$\mathcal{P}(X, Y) = \bigoplus \mathcal{P}_i(X, Y)$ \ill{}
\uncertain{Soit} \ill{}}
\marginal{\ill{} $\mathcal{P}($\ill{}$)$ \ill{} \ill{}
($n_i m_j \neq 0$) \ill{} \ill{}}
\note{deux notes de biais : la première à droite, en face de e), la
seconde dans le coin inférieur gauche, en face de f) ; l'une et l'autre
presque entièrement illisibles. Les axiomes vont ici jusqu'à g), alors
que la page 83 cite « les axiomes a) à f) ».}

Si $k$ est un anneau commutatif et si $\mathcal{P}_1$ est munie d'une
structure $k$-linéaire par
\[
k \longrightarrow \mathrm{End}(\mathcal{P}_1) \longleftarrow
\mathrm{End}(\mathcal{P})
\]
\struck{\ill{} \ill{} que} telle que les \uncertain{relations} plus haut
sont vraies pour $\lambda \in k^{r}$ (non seulement $\lambda \in
\mathbb{Z}^{r}$), on dit que $\mathcal{P}$ est $k$-polynomiale.

\page{77}
\subsection*{Relations avec produits $\otimes$ et $\Gamma^{p}$}
\note{titre souligné sur la page}

La donnée d'une cat.\ polynomiale $\mathcal{P}$ permet de définir un
foncteur \uncertain{trimultiplicatif}
\[
\mathcal{P}_1^{\circ} \times \mathcal{P}_1^{\circ} \times \mathcal{P}_1
\longrightarrow \mathit{Ab}, \qquad
(X, Y; Z) \longmapsto \mathcal{P}_{(1,1)}(X, Y; Z)
\]
(et de \uncertain{même} pour plus d'arguments $2$ \uncertain{facteurs})
Si, pour $X, Y$ fixés, $Z$ variable, ce foncteur est représentable,
l'objet qui le représente est noté $X \otimes Y$. Il dépend
fonctoriellement de $X, Y$ dans $\mathcal{P}_1$ (mais non pour $X, Y$
\struck{\ill{}} considérés comme dans $\mathcal{P}$\,!).
\struck{\ill{}} On voit que $X \otimes Y$ existe ssi $Y \otimes X$
existe, et \uncertain{ils} \uncertain{sont} \uncertain{canon.}\
isomorphes (\uncertain{par} $s_{X,Y}$), \uncertain{avec}
$s_{X \otimes Y} \circ s_{Y \otimes X} = \mathrm{id}_{Y \otimes X}$.

Soit $X$, $n \in \mathbb{N}$ fixés, si le foncteur
\[
Y \longmapsto \mathcal{P}_n(X, Y)
\]
sur $\mathcal{P}_1$ est représentable, l'objet qui le représente est noté
$\Gamma^{n} X$ ; il \uncertain{dépend} encore fonctoriellement de $X$ dans
$\mathcal{P}_1$. L'objet $\Gamma^{0}(X)$ existe ssi $\Gamma^{0}(0)$
existe, il est noté $\mathbf{1}$.

Prouvons que \struck{\ill{}} si $\mathbf{1}$ existe, alors
$\forall\, X \in \mathrm{Ob}\,\mathcal{P}$,
\[
\mathbf{1} \otimes X \simeq X \quad \text{i.e.} \quad
\mathcal{P}_{(1,1)}(\mathbf{1}, X; Y) \simeq \mathcal{P}_1(X, Y)
\; (= \mathrm{Hom}_{\mathcal{P}_1}(X, Y)) \ ?
\]
\note{sous le signe $\simeq$, « iso bifonctoriel »}
et plus particulièrement $\mathbf{1} \otimes \mathbf{1} \simeq
\mathbf{1}$ i.e.\ $\forall\, Y$
\[
\mathcal{P}_{1,1}(\mathbf{1}, \mathbf{1}; Y) \simeq
\mathcal{P}_1(\mathbf{1}, Y) \ ?
\]

\page{78}
On est obligé de renforcer la propriété universelle de $\otimes$ et
$\Gamma^{p}$, pour avoir les propriétés d'assoc.\ etc.\ qu'on
\uncertain{souhaite}.

Soit $X \times Y \to$ \struck{\ill{}} $T$ un morphisme de bidegré
$(1,1)$, alors pour tt $Z$ on a
\[
X \times Y \times Z \longrightarrow T \times Z
\]
d'où pour tt $M$
\[
\underset{1,n}{\mathcal{P}}(T, Z; M) \longrightarrow
\mathcal{P}_{1,1,n}(X, Y, Z; M)
\]
On dit que $T$ est un \emph{produit tensoriel strict} de $X$ et $Y$ si
c'est un \uncertain{bijection} \struck{\ill{} $X \otimes Y$ \ill{}}
pour tt $n$ et tout $M$. Faisant $n = 0$, $Z = 0$, on trouve
\[
\mathcal{P}_{1}(T; M) \simeq \mathcal{P}_{1,1}(X, Y; M)
\]
\note{l'indice de $\mathcal{P}$ dans le membre de gauche est surchargé ;
« $Z = 0$ » est écrit au-dessus de la ligne}
donc $T$ est un produit tensoriel. C'est
\uncertain{dire} \uncertain{qu'on} \uncertain{exigera} que \ill{} des
$\otimes$ \uncertain{strict} \struck{\ill{}} \ill{} \ill{} \ill{}
\uncertain{la} \uncertain{définition} \uncertain{précédente}
\struck{\ill{} \ill{} \ill{}} \ill{} \ill{} soit
\uncertain{associative} \ill{} \uncertain{compatible} avec
\[
\mathcal{P}_1(X_1 \otimes \cdots \otimes X_r; M) \simeq
\mathcal{P}_{1, \ldots, 1}(X_1, \ldots, X_r; M)
\]
\note{un $\mathrm{Hom}_{\mathcal{P}_1}$ biffé en tête de la formule}
\note{les deux lignes qui précèdent la formule sont surchargées
d'interlignes, en partie biffés ; on n'en donne que ce qui se lit}
On sait qu'un objet unité (sans \uncertain{référence} d'une donnée
d'associativité) s'il existe, est unique. Il faut examiner pourquoi ce
devrait être $\Gamma^{0}(0)$.
\note{un point d'interrogation dans la marge gauche, en face de cette
phrase}

Revenons aux $\Gamma^{n}$, soit $X = X_1 \times \cdots \times X_r$, tel
que $\Gamma^{n}(X)$ existe. Considérons les \struck{\ill{}} foncteurs
$\mathcal{P}_{\nu = (\nu_1, \ldots, \nu_r)}(X_1, \ldots, X_r; M)$
\add{$\mathcal{P}_1 \to (\mathit{Ab})$}, les foncteurs (\uncertain{cov.})
représentés par $\Gamma^{n}(X)$, supposons qu'ils soient aussi
\uncertain{repr.}\ (automatiques si $\mathcal{P}_1$ est
karoubienne) \uncertain{par} $\Gamma^{\nu_1, \ldots, \nu_r}(X_1, \ldots,
X_r) = \Gamma^{\nu}(X_1, \ldots, X_r)$ l'\add{objet} \struck{\ill{}} qui
le représente. \struck{\ill{}} On va définir (sans \uncertain{réserve}
d'existence des $\Gamma^{\nu_i} X_i$ \add{et des} produits tensoriels,
\ldots\ ou dans $\check{\mathcal{P}}_1 =
\widehat{\mathcal{P}_1^{\circ}}^{\circ}$ si \ldots
\note{la page s'arrête sur « si » ; la suite est à la page 79}

\page{79}
on veut \uncertain{faire} abstraction des questions de représentabilité
\uncertain{sauf} pour les $\Gamma^{\nu_i}(X_i)$)
\[
\Gamma^{\nu_1}(X_1) \otimes \cdots \otimes \Gamma^{\nu_r}(X_r)
\longleftarrow \Gamma^{\nu}(X_1, \ldots, X_r)
\]
\note{le dernier argument est écrit $\alpha_r$ (ou $\mathcal{X}_r$) sur la
page}
en \uncertain{donnant} un hom fonctoriel en $M$ dans $\mathcal{P}_1$
\[
\mathcal{P}^{\mathfrak{g}}_{\nu}(X_1, \ldots, X_r; M) \longleftarrow
\mathcal{P}_{1, \ldots, 1}(\Gamma^{\nu_1} X_1, \ldots,
\Gamma^{\nu_r} X_r; M)
\]
\note{l'exposant de $\mathcal{P}$, lu $\mathfrak{g}$, est d'une lecture
douteuse}
par composition avec les morphismes canoniques de $\mathcal{P}$
\[
X_1 \xrightarrow{\ \gamma_{\nu_1}\ } \Gamma^{\nu_1} X_1, \quad \ldots,
\quad X_r \xrightarrow{\ \gamma_{\nu_r}\ } \Gamma^{\nu_r}(X_r)
\]
et on suppose que c'est \add{\uncertain{toujours}} un iso. Moyennant
quoi, on a
\[
\Gamma^{n}(X_1 \times \cdots \times X_r) \xrightarrow{\ \sim\ }
\sum_{\substack{\nu = (\nu_i) \in \mathbb{N}^{r} \\ \sum \nu_i = n}}
\underbrace{\Gamma^{\nu_1}(X_1) \otimes \cdots \otimes
\Gamma^{\nu_r}(X_r)}_{\Gamma^{\nu_1 \ldots \nu_r}(X_1, \ldots, X_r)}
\]
[\emph{NB} Il suffirait de l'établir \ldots\ \uncertain{avec} les
hypothèses \uncertain{réduites} $r = 2$ --- \uncertain{puis} $r = 2$
\uncertain{d'abord} et la déduire \uncertain{ensuite} pour $r$
\uncertain{quelc.}

En particulier, on trouve, \struck{pour} \uncertain{si} \ill{}
considère le produit $X \times 0$, pour $\mathbf{1} = \Gamma^{0}(0)$
\[
X \otimes \mathbf{1} = \Gamma^{1}(X) \otimes \Gamma^{0}(0) \simeq
\Gamma^{1,0}(X, 0) \simeq X
\]
car $\mathcal{P}^{\mathfrak{g}}_{1,0}(X, 0; M) \simeq \mathcal{P}_1(X, M)$,
et de \uncertain{même} $\mathbf{1} \otimes X \simeq X$
\struck{\ill{}}.
\note{au-dessus de la formule, une première écriture biffée, « $\Gamma^{1}(X \times 0) \simeq$ » ; sous la formule, un « $X \otimes \mathbf{1}$ » biffé}

On suppose qu'il y a des produits tensoriels stricts \uncertain{dans}
\uncertain{les} \uncertain{objets} \uncertain{de} $\mathcal{P}_1$, et que
les $\Gamma^{n} X$ satisfont à la condition plus haut.
\marginal{\ill{} \ill{} \ill{} \ill{}}
\note{un trait vertical dans la marge gauche, en face de cette phrase, et
une note de biais de quelques mots, illisible ; le bas du feuillet est
blanc}

\page{80}
Ainsi
\[
X \longmapsto \Gamma^{*}(X) = (\Gamma^{i}(X))_{i \geq 0}
\]
est un foncteur (\uncertain{pas} additif\,!)
\[
\mathcal{P}_1 \longrightarrow \mathrm{Grad}(\mathcal{P}_1)
\]
qui est en fait un foncteur « transformant sommes en $\otimes$ ». Comme
tout objet de $\mathcal{P}_1$ est un \add{objet} groupe commutatif pour
l'opération $\oplus$, il s'ensuit que sur chaque $\Gamma^{*}(X)$ il y a une
structure de $\otimes$-algèbre \add{graduée} commutative, fonctorielle en
$X$.
\marginal{\uncertain{Pour} \ill{} \ill{} \ill{} \ill{} \ill{}
$\otimes$-algèbres \ill{} \ill{} \ill{} \ill{}}
\note{note de biais dans la marge gauche, en face de ce passage et du
suivant, presque entièrement illisible ; un « 2 » au-dessous}

\struck{Explicitation} des $\mathcal{P}_1$
\note{« des $\mathcal{P}_1$ » est écrit au-dessus du mot biffé}
On se donne des homs, fonctoriels en $X$
\[
(*) \qquad \Gamma^{p}(\Gamma^{q}(X)) \longleftarrow \Gamma^{pq}(X)
\]
\note{le signe devant la formule, lu $(*)$, est surchargé}
ou, ce qui revient au même, un hom de degré $pq$
\[
X \longrightarrow \Gamma^{p}(\Gamma^{q}(X))
\]
en prenant donc le composé
\[
X \xrightarrow[\text{hom.\ univ.\ de degré } q]{\ e^{X}_{q}\ }
\Gamma^{q}(X)
\xrightarrow[\text{hom.\ univ.\ de degré } p]{\ e^{\Gamma^{q}(X)}_{p}\ }
\Gamma^{p}(\Gamma^{q}(X))
\]
On a la condition « d'associativité »

\begin{tikzcd}[column sep=huge, row sep=large]
\Gamma^{pqr}(X) \arrow[r, "\alpha^{X}_{p,qr}"] \arrow[d, "\alpha^{X}_{pq,r}"'] & \Gamma^{p}(\Gamma^{qr}(X)) \arrow[d, "{\Gamma^{p}(\alpha^{\Gamma^{r}(X)}_{q,r})}"] \\
\Gamma^{pq}(\Gamma^{r}(X)) \arrow[r, "\alpha^{\Gamma^{r}(X)}_{p,q}"'] & \Gamma^{p}(\Gamma^{q}(\Gamma^{r}(X)))
\end{tikzcd}

\note{le diagramme est marqué $(*)$ à gauche et « comm. » au centre. Les
indices et exposants des $\alpha$ sont retouchés ; on les donne comme ils
se lisent, en particulier l'exposant $\Gamma^{r}(X)$ de la flèche verticale
de droite, où le contexte attendrait $\alpha^{X}_{q,r}$.}
et d'unité : \uncertain{l'isom.}\ \ill{} \ill{} (\ill{} \ill{})
\uncertain{\ill{}} \uncertain{pour}
\[
\Gamma^{1} X = X, \qquad \Gamma^{0} X = \mathbf{1}
\]
\note{les conditions d'unité sont énoncées en détail, comme (**) et
(***), en tête de la page 81}

\end{document}
